% 学习报告(随项目每周迭代) —— 编译方式:XeLaTeX
% 结构:第一部分 = 项目启动前的初版报告(2026-07-14 docx 原文转录,不改动);
%       第二部分 = 项目进行中的每周迭代,只增不删,并逐步回应第一部分第 7 节的疑问
\documentclass[11pt]{ctexart}
\usepackage[a4paper, margin=2.5cm]{geometry}
\usepackage{amsmath, amssymb}
\usepackage{booktabs}
\usepackage[hidelinks]{hyperref}
\usepackage{enumitem}
% generated by analysis/week12/04_report_numbers.py - do not edit by hand
\newcommand{\rnWnineQuoteMs}{1.8}
\newcommand{\rnWnineRedisMs}{0.15}
\newcommand{\rnWnineRenderMs}{724}
\newcommand{\rnWnineDailyS}{4.5}
\newcommand{\rnMemIntervals}{670}
\newcommand{\rnMemTickers}{666}
\newcommand{\rnPanelTenRows}{617{,}211}
\newcommand{\rnPanelTenTickers}{532}
\newcommand{\rnEone}{9.5}
\newcommand{\rnEtwo}{6.8}
\newcommand{\rnEthree}{5.4}
\newcommand{\rnEfour}{163.8}
\newcommand{\rnEfive}{69.9}
\newcommand{\rnTenDailyS}{38.0}
\newcommand{\rnGoneDWin}{72}
\newcommand{\rnGoneDPooled}{0.812}
\newcommand{\rnRoundOneRidge}{7.16}
\newcommand{\rnRoundOneGbdt}{2.23}
\newcommand{\rnNoPrices}{51}
\newcommand{\rnInStudy}{597}
\newcommand{\rnTruncated}{3}
\newcommand{\rnPostExit}{7}
\newcommand{\rnLatestBeforeMs}{457}
\newcommand{\rnLatestBeforeRows}{579}
\newcommand{\rnLatestFinalMs}{0.3}
\newcommand{\rnLatestFinalRows}{1496}
\newcommand{\rnQuBeforeMs}{390}
\newcommand{\rnQuBeforeN}{578}
\newcommand{\rnQuAfterMs}{23.5}
\newcommand{\rnQuAfterN}{1496}
\newcommand{\rnMarketBeforeS}{12.8}
\newcommand{\rnMarketAfterS}{2.6}
\newcommand{\rnFeatBeforeS}{14.7}
\newcommand{\rnFeatAfterS}{14.5}
\newcommand{\rnDailyAfterS}{25.8}
\newcommand{\rnScaleSteadyS}{59.5}
\newcommand{\rnScaleN}{1517}
\newcommand{\rnScaleBackfillS}{530}
\newcommand{\rnRenderBeforeMs}{591}
\newcommand{\rnRenderAfterMs}{514}
\newcommand{\rnFinalQuMs}{50}
\newcommand{\rnFinalQuN}{1496}
\newcommand{\rnFinalScalarMs}{27}
\newcommand{\rnFinalXsMs}{27}
\newcommand{\rnFinalRenderMs}{510}
\newcommand{\rnGarchOrdS}{3.2}
\newcommand{\rnGarchRefitS}{15.2}
\newcommand{\rnGarchRefits}{1{,}326}
\newcommand{\rnGarchTickers}{1{,}479}
\newcommand{\rnCatchupS}{84.4}
\newcommand{\rnCatchupFeatS}{62.7}
\newcommand{\rnDbFinalMb}{1601}
\newcommand{\rnDbElevenMb}{1570}
\newcommand{\rnRssFinalMb}{299}
\newcommand{\rnRssElevenMb}{286}
\newcommand{\rnElevenTestMl}{0.755}
\newcommand{\rnElevenTestG}{0.791}
\newcommand{\rnElevenMaeMl}{7.91}
\newcommand{\rnElevenMaeG}{8.36}
\newcommand{\rnElevenMaeE}{10.35}
\newcommand{\rnElevenVal}{0.705}
\newcommand{\rnElevenTestStart}{2024-05-24}
\newcommand{\rnNoDummyTest}{0.741}
\newcommand{\rnUnknownN}{52}
\newcommand{\rnUnknownRemoved}{51}
\newcommand{\rnGTrain}{0.840}
\newcommand{\rnGVal}{0.777}
\newcommand{\rnCsOnly}{0.742}
\newcommand{\rnLevelOnly}{0.801}
\newcommand{\rnOosStart}{2022-12-27}
\newcommand{\rnOosEnd}{2024-12-31}
\newcommand{\rnOosDates}{506}
\newcommand{\rnOosRows}{244{,}038}
\newcommand{\rnOosTickers}{505}
\newcommand{\rnOosMaeE}{7.55}
\newcommand{\rnOosMaeG}{6.49}
\newcommand{\rnOosMaeMl}{6.15}
\newcommand{\rnHOneDelta}{1.07}
\newcommand{\rnHOneLo}{0.65}
\newcommand{\rnHOneHi}{1.62}
\newcommand{\rnHOneSizeTicker}{48}
\newcommand{\rnHOneSizeSector}{9.3}
\newcommand{\rnHOnePholm}{0.002}
\newcommand{\rnHTwoDelta}{0.337}
\newcommand{\rnHTwoLo}{0.206}
\newcommand{\rnHTwoHi}{0.456}
\newcommand{\rnHTwoSizeTicker}{50}
\newcommand{\rnHTwoSizeSector}{4.7}
\newcommand{\rnHTwoPholm}{0.002}
\newcommand{\rnHThreeDelta}{1.48}
\newcommand{\rnHThreeLo}{1.15}
\newcommand{\rnHThreeHi}{1.80}
\newcommand{\rnHThreeSizeTicker}{23}
\newcommand{\rnHThreeSizeSector}{3.7}
\newcommand{\rnHThreePholm}{0.002}
\newcommand{\rnHFourDelta}{1.77}
\newcommand{\rnHFourLo}{1.29}
\newcommand{\rnHFourHi}{2.44}
\newcommand{\rnHFourSizeTicker}{42}
\newcommand{\rnHFourSizeSector}{8.0}
\newcommand{\rnHFourPholm}{0.002}
\newcommand{\rnHTwoYThree}{0.048}
\newcommand{\rnHTwoYFour}{0.634}
\newcommand{\rnHTwoBDelta}{0.286}
\newcommand{\rnSelfIidTicker}{6.7}
\newcommand{\rnSelfShockTicker}{69}
\newcommand{\rnSelfShockSector}{9.3}
\newcommand{\rnIvJpm}{0.710}
\newcommand{\rnIvJpmCorr}{0.78}
\newcommand{\rnIvJpmCorrConst}{0.83}
\newcommand{\rnIvPanelRows}{467{,}189}
\newcommand{\rnIvPanelTickers}{506}
\newcommand{\rnIvMaeIv}{5.16}
\newcommand{\rnIvMaeGd}{6.63}
\newcommand{\rnIvMaeGy}{6.92}
\newcommand{\rnIvMaeRv}{7.60}
\newcommand{\rnIvMaeEw}{7.52}
\newcommand{\rnIvWinGd}{92}
\newcommand{\rnIvWinGy}{70}
\newcommand{\rnIvConstIv}{4.69}
\newcommand{\rnIvConstGy}{6.47}
\newcommand{\rnIvXsIv}{0.95}
\newcommand{\rnIvXsGy}{0.84}
\newcommand{\rnIvTsIv}{0.11}
\newcommand{\rnIvGapActual}{1.77}
\newcommand{\rnIvGapConst}{1.78}
\newcommand{\rnBandGCalm}{5.48}
\newcommand{\rnBandGFear}{5.86}
\newcommand{\rnBandESameCalm}{6.66}
\newcommand{\rnBandESameFear}{4.89}
\newcommand{\rnBandEFullCalm}{9.01}
\newcommand{\rnBandEFullFear}{10.83}
\newcommand{\rnBandGDays}{1011}
\newcommand{\rnToolAsOf}{2026-10-02}
\newcommand{\rnToolQuoted}{1{,}493}
\newcommand{\rnToolGarch}{1{,}478}
\newcommand{\rnToolEwma}{15}
\newcommand{\rnToolMl}{496}

\setlist{nosep}

\title{\textbf{期权定价与时序建模学习报告}\\
\large (初版完成于项目启动前,随项目进度每周迭代更新)}
\author{佟铸城 (Zhucheng Tong)}
\date{初版:2026 年 7 月 14 日 \quad 最近更新:2026 年 10 月 13 日(截至第 12 周,结项)}

\begin{document}
\pagenumbering{roman}   % title page + table of contents are not counted; the body starts at page 1
\maketitle
\tableofcontents
\clearpage
\pagenumbering{arabic}

\part*{第一部分\quad 项目启动前的预备学习(2026-07-14 初版)}
\addcontentsline{toc}{part}{第一部分\quad 项目启动前的预备学习}

\noindent 本次学习材料:
\begin{itemize}
  \item Hull《Options, Futures, and Other Derivatives》第 11 版(第 13、15、26 章为主);
  \item 哥大课程讲义 STAT5265《Stochastic Methods in Finance》(Unit 1、3)、STAT5263《Statistical Inference and Time Series Modelling》(第 4、7、9 章);
  \item 《动手学深度学习》(\url{https://zh.d2l.ai})中的 RNN 与 Transformer 章节;
  \item 南瓜书第 5 章作为神经网络推导参考。
\end{itemize}

\section{我对本次任务的理解}
我拿到``奇异期权、基础定价模型、深度学习时序模型''三个主题时首先想到的是它们为什么会被放在一起,而读完材料中相关的部分后,我的理解是这三者其实是同一条线上的三个不同节点——期权定价的本质是对标的资产动态的建模,一些经典模型(二叉树、BSM)在特定的强假设下给出解析框架;奇异期权因为收益结构复杂而会最先暴露出上述解析方法的边界,让我们需要借助数值方法;而经典模型里最脆弱的常数波动率假设又刚好是时序模型(包括从 GARCH 到深度学习)要解决的问题。本报告将按这条线组织和记录我的学习过程和形成的理解。

\section{基础期权定价模型}

\subsection{二叉树:定价公式里为什么没有``预测''}
Hull 第 13 章的单步二叉树最引起我注意的一点是,标的资产的真实上涨概率完全不出现在期权价格里——从风险中性概率
\[ p = \frac{e^{r\Delta t} - d}{u - d} \]
的公式看,它只由无风险利率和涨跌幅度决定。我对此的理解是,定价的出发点不是``预测会涨还是会跌'',而是构造一个股票+期权的组合,以此来使它在涨、跌两种情形下价值相同。方向性风险被对冲掉之后,组合只能赚无风险收益,期权价格由此被无套利条件锁死。所谓``风险中性概率''也并不是任何人对市场的看法,而是把这个无套利价格改写成期望形式时自然出现的权重。这一点让我重新理解了衍生品定价和股票投资的根本区别:前者做的是复制与对冲,后者才需要预测方向。

\subsection{BSM:PDE 路线与测度变换路线的对照}
Hull 第 15 章和 STAT5265 Unit 1 分别从两个方向推导了同一个公式,对照着读收获很大:第一个是 Hull 的路线(PDE):构造 delta 对冲组合,用 Ito 引理展开后发现 $dz$ 项被消掉,组合瞬时无风险,从而得到 BSM 偏微分方程,真实漂移率 $\mu$ 在这一步物理上地消失了。第二个是讲义的路线(鞅/Girsanov):组合不改动,只改概率的测度,通过书中 Girsanov 定理把布朗运动的漂移换掉,使贴现后的股价在新测度 $Q$ 下成为鞅,期权价格就是 $Q$ 下期望收益的贴现。$\mu$ 在换测度时被吸收进了新的布朗运动。看起来,这两条路线做的是同一件事,那就是消灭 $\mu$,不过一个在对冲组合层面消,一个在概率测度层面消(后者在数学上更通用,比如想推广到随机利率、多资产时几乎只能走测度路线)。二叉树里``真实概率不出现''和连续时间里``$\mu$ 不出现''分别是同一个事实在离散与连续两种情况下的表述。

\subsection{哪条假设最先失效}
BSM 的假设里,我认为现实中最先失效、也最重要的一条是常数波动率:同一标的、同一到期日、不同行权价的期权反解出的隐含波动率并不是一条水平线,而会是微笑或偏斜的曲线,这意味着 $\sigma$ 在这些情形下不会是常数。其次是连续对冲(现实中有交易成本和离散调仓)和无跳跃(现实中存在财报、政策事件)。常数波动率的失效为第四节埋下伏笔:如果 $\sigma$ 本身是随时间变化、需要建模的量,那么波动率的时序建模就成为定价体系的一部分。

\section{奇异期权}

\subsection{我的分类方式}
Hull 在第 26 章连续列举了十几种奇异期权。为了避免被绕晕,我按它们对路径的依赖程度把它们重新归了类,因为我发现这个维度直接决定定价方法的选择:
\begin{enumerate}
  \item \textbf{欧式改造型(只看终值)}:如 gap、binary/digital、chooser——这些基本可沿用 BSM 框架改边界条件;
  \item \textbf{路径依赖型}:如 barrier、Asian、lookback——这些需要知道整条路径的分布信息;
  \item \textbf{多资产型}:如 basket、rainbow、exchange——这些的难点在维度灾难和相关性结构。
\end{enumerate}

\subsection{一个很有助于我理解的对比:障碍有闭式解,亚式没有}
Hull 对障碍期权直接给了公式,没有解释来源;STAT5265 Unit 3 恰好补上了推导:障碍期权之所以有闭式解、是因为布朗运动的反射原理给出了``终值+运行最大值''的联合分布,路径依赖信息(是否触碰过障碍)被压缩成一个分布已知的统计量。对照之下,算术平均亚式期权没有闭式解的原因也清楚了:对数正态变量之和的分布不可解。同样是路径依赖,``最大值''这个统计量的分布可解,``算术平均''不可解,所以是否存在闭式解不取决于期权看起来多复杂,而取决于其收益所依赖的路径统计量是否有可解的分布。还有个有趣的旁证:把算术平均换成几何平均(对数正态的几何平均仍是对数正态),闭式解就会立刻恢复,这是后面 6.2 节实验里控制变量法的理论基础。

\subsection{数值方法的选择逻辑}
综合 Hull 第 21、26 章和我此前的项目经验,我把数值方法的适用边界整理为:树和有限差分适合低维问题且天然支持美式提前行权(可以在每个节点比较行权与持有价值);Monte Carlo 适合高维和路径依赖问题(收敛速度与维数无关),但处理提前行权困难,因为模拟是沿时间正向的,而行权决策需要知道``继续持有的价值''这一逆向信息。这里可以结合我此前做 CEV-Heston 条件 Monte Carlo 的经验:MC 的原始收敛速度是 $1/\sqrt{n}$ 量级,很慢,各种方差缩减技术(条件化、控制变量)在实操中不仅是锦上添花、而是必需品——当时是通过条件化把长期限($T=10$)下的定价误差从基线的 19.8 个标准误压到 2 个标准误以内,而本报告 6.2 节则用一个更简单的例子(几何亚式作控制变量)复现了这一逻辑。

\subsection{市场为什么需要奇异期权}
从业务角度想了一下``为什么会存在这些产品'':一是定制化风险敞口——客户的风险本身就是路径依赖的(例如持续采购原材料的企业,其成本暴露天然是``平均价格'',亚式期权是精确匹配而非近似);二是降低权利金——敲出障碍期权比同参数 vanilla 便宜,因为买方放弃了一部分情景下的收益,本质是让客户只为自己真正需要的保护付费。

\section{从经典时序模型到深度学习}

\subsection{经典基线:ARMA 建模均值,GARCH 建模波动率}
重读 STAT5263 时我注意到一个以前没有在意的分工:ARMA(第 4 章)建模的是序列的条件均值,而 ARCH/GARCH(第 7 章)建模的是条件方差,对期权业务而言重要的是后者——2.3 节说过,BSM 最脆弱的假设是常数波动率,而讲义第 7 章用 GARCH 对 S\&P500 建模的例子正说明波动率有很强的可预测结构(聚集性、均值回复)。期权定价关心的时序问题,主线是波动率而非价格方向——这和 2.1 节``定价不需要预测方向''的结论互相印证,我认为这是本次学习中三个主题之间相当重要的一条连接线。

\subsection{一条演进线:状态空间 $\to$ RNN $\to$ LSTM $\to$ Transformer}
STAT5263 第 9 章的状态空间模型给了我理解 RNN 的一个桥梁。状态空间模型用隐状态概括历史、用状态方程 $x_t = F x_{t-1} + w_t$ 演化;对照 d2l 中 RNN 的隐状态更新 $h_t = \tanh(W_h h_{t-1} + W_x z_t)$,这两者的结构几乎是一一对应的——RNN 可以看作非线性、参数由数据学出来的状态空间模型。沿着``每个模型被用于解决前一个的什么问题''往下走大概能得到如下的思路:RNN:能用隐状态携带任意长的历史信息,摆脱了前面 ARMA 的线性和固定阶数限制;但反向传播沿时间连乘会导致梯度消失,实际上记不住长程依赖。LSTM/GRU:引入了门控,让``记住/遗忘''变成可学习的决策,细胞状态提供了一条梯度可以近乎无衰减流过的通道,解决长程依赖。Transformer:干脆去掉循环,用自注意力让任意两个时点直接交互,由此任意距离的依赖都是``一步''可达,且可以并行训练;代价是失去天然的顺序概念,需要靠位置编码补回。

\subsection{我的保留意见:DL 在期权业务里的合理落点}
实际上,我个人对``用深度学习直接预测价格''持一定的怀疑态度:金融时序信噪比极低、有效样本少(长历史不等于同分布样本)、且非平稳,总感觉大容量模型在这种数据上最容易学到的是噪声。我认为 DL 更可靠的落点是那些信号结构更强的环节:波动率及波动率曲面的建模与拟合(结构性强、数据充足)、定价加速(用网络作复杂模型的 surrogate,训练数据由模型自身生成、要多少有多少)、以及对冲策略学习(目标是最小化对冲误差而非预测点位)。这些方向共同点是:不与``市场近似有效''这一事实正面冲突,不过这个判断有待后续的检验。

\section{Python 环境检查}
按要求检查了本机环境(macOS,Python 3.13.5),核心库应该没有缺的,所有实验都是以 .py 脚本运行、图表输出为 PNG 文件:numpy 2.2.6(数值计算)、scipy 1.16.1(正态分布、优化)、pandas 2.3.1(数据处理)、matplotlib 3.10.5(绘图)、scikit-learn 1.7.1(经典 ML)、statsmodels 0.14.5(ARIMA/GARCH 基线)、torch 2.8.0(深度学习,时序模型部分)。

\section{两个验证性小实验}
为了检验第二、三节的理解(也顺带验证环境可用),写了两个小脚本。

\subsection{二叉树收敛到 BSM(binomial\_convergence.py)}
参数 $S_0=K=100$,$r=3\%$,$\sigma=20\%$,$T=1$。BSM 闭式解为 9.4134;CRR 树 $N=10$ 时价格 9.2179(误差 0.196),$N=100$ 时 9.3936(误差 0.020),$N=500$ 时 9.4094(误差 0.004)。从图里看到两点:(1) 误差在 log-log 图上贴着 $1/N$ 参考线下降,验证了离散时间模型与连续时间模型在极限下的一致性——由此 2.1 与 2.2 节``同一事实的离散/连续两种表述''有了数值证据;(2) 收敛是振荡的,价格围绕 BSM 上下摆动,原因是终端节点与行权价的相对位置随 $N$ 的奇偶变化——这表示实操中取相邻两个 $N$ 的平均可以显著加速收敛。

\subsection{Monte Carlo 定价算术亚式期权 + 控制变量(mc\_asian\_option.py)}
算术平均亚式看涨(月度观测 12 次,同上参数)无闭式解,用 10 万条路径的 MC 定价;并以几何平均亚式(3.2 节:其闭式解存在,算出为 5.4353)作控制变量。结果:普通 MC:价格 5.6346,标准误 0.0261;MC+控制变量:价格 5.6326,标准误 0.0007。两个观察:(1) 算术与几何亚式收益的相关系数高达 0.9996,控制变量把标准误压缩了 36 倍——等效于普通 MC 用 1300 多倍的路径数,印证了 3.3 节``方差缩减是必需品''的说法;(2) 算术亚式价格(约 5.63)显著低于同参数欧式期权(9.41),因为平均化压低了收益的波动——这也解释了 3.4 节里亚式期权对冲``平均成本''型风险时更便宜、更精确的业务逻辑。

\section{疑问与下一步}\label{sec:questions}
学习过程中产生了几个还没有答案的问题,留作可能的后续方向:
\begin{enumerate}
  \item \textbf{美式/百慕大型奇异期权的 MC 定价}:3.3 节提到 MC 难以处理提前行权,本人查到 Longstaff--Schwartz 最小二乘法是标准解法,但还没有弄清楚回归基函数的选择对精度的影响有多大。
  \item \textbf{波动率微笑与奇异期权定价的一致性}:既然微笑曲线证伪了常数波动率,实务中给障碍期权定价时用哪个波动率?是否必须升级到局部波动率或随机波动率模型才能自洽?
  \item \textbf{DL 定价模型的无套利约束}:如果用神经网络拟合定价函数或波动率曲面,如何才能保证输出满足无套利条件(如价格对行权价的凸性)?是另加软约束、还是有结构性的做法?
  \item \textbf{Transformer 在低信噪比金融时序上是否真的优于 LSTM}:d2l 的例子都是文本类高信号数据,金融数据上注意力机制的优势是否还成立仍需要找实证文献确认。
\end{enumerate}

% ====================================================================
\part*{第二部分\quad 项目进行中的每周迭代}
\addcontentsline{toc}{part}{第二部分\quad 项目进行中的每周迭代}

\section{第 1 周:Chooser Option 与数据工程}

\subsection{Chooser option 的解析定价(Rubinstein, 1991)}
项目标的正是初版 3.1 节分类中``欧式改造型''里的 chooser option,本周把它的解析解推导补全了。Chooser option 赋予持有人在\textbf{选择日} $t_1$ 决定期权为欧式 call 或 put 的权利,两者共享行权价 $K$ 与到期日 $T_2$($t_1 < T_2$),选择日价值为
\[
V(t_1) = \max\big( C(S_{t_1}, K, T_2 - t_1),\; P(S_{t_1}, K, T_2 - t_1) \big).
\]
利用 put--call parity($P = C - S e^{-q\tau} + K e^{-r\tau}$,$\tau = T_2 - t_1$)可改写为
\[
V(t_1) = C + \max\big( 0,\; K e^{-r(T_2 - t_1)} - S_{t_1} e^{-q(T_2-t_1)} \big),
\]
即 simple chooser 等价于\textbf{静态组合}:一份 $(K, T_2)$ 的 call,加一份 $(K e^{-r(T_2-t_1)}, t_1)$ 的 put(标的经股息调整)。因此 BSM 假设下 chooser 有闭式解——这再次印证初版 3.2 节的判断:有没有闭式解,取决于收益所依赖的统计量分布是否可解;chooser 的``路径依赖''被 parity 压缩成了单时点量。直觉上 chooser 比 straddle 便宜:选择在 $t_1$ 锁定,straddle 相当于把选择拖到 $T_2$。

\subsection{参考论文精读与复现规划}
精读了项目参考论文(Huang, Wang \& Wan, 2021,复现参数:$S_0=156.7$、$K=150$、$r=0.15\%$、$\sigma=28.2\%$、$q=2.33\%$、$t_1=0.5$、$T_2=1$)。论文用两段 GBM 模拟定价,但其选择规则是在 $t_1$ 比较 $S_{t_1}$ 与 $K$(moneyness),而非比较两个期权的价值 $\max(C,P)$——由 8.1 节的分解可知理性临界点并不在 $K$ 处。第 3 周复现时将同时实现两种规则,量化``天真规则''的价值低估幅度;论文未给出闭式解基准价,复现将以敏感性曲线形状对照 + 闭式解/MC 收敛验证为准(与导师指定的蒙特卡洛验证思路一致,也与初版 3.3 节``MC 是路径依赖问题的通用武器''相衔接)。

\subsection{数据工程实践}
本周完成数据管道搭建(2018--2024 日频:JPM、VIX、标普 500、国债利率),几条实操教训:(1) \textbf{免费数据源的能力边界需实测}:Alpha Vantage 免费档的全量历史接口已转付费(仅返最近约 100 个交易日),文档与实际行为不同步,架构上已将其降级为备用源、主源改用 yfinance;(2) \textbf{多源冗余}:国债收益率同时取 FRED 与 Yahoo 互为校验;(3) \textbf{日历异质性}:宏观数据(工作日日历,1850 行)与市场数据(NYSE 交易日历,1760 行)需显式对齐而非默认合并;(4) 工程规范:API key 经 \texttt{.env} 管理并排除出版本控制,采集器模块化以便接入自动化调度。

\subsection{对第一部分疑问的回应进度}
初版第 \ref{sec:questions} 节的四个疑问,预计的回应时点:疑问 2(波动率微笑与定价一致性)与疑问 3(DL 定价的无套利约束)将在第 5--6 周 ML 建模时正面遭遇;疑问 4(Transformer vs LSTM)第 5--6 周的模型对比可提供本项目数据上的实证;疑问 1(LSM 定价美式型)超出本项目核心范围,若进入 enhancement 阶段(美式 chooser 是自然的扩展方向)再处理。

% ====================================================================
% 以下章节随项目进度逐周补充
% ====================================================================

\section{第 2 周:数据预处理与特征工程}

\subsection{缺失值填充要匹配数据生成机制}
本周清洗中体会最深的一点:填充方法的选择不是技术偏好,而是对\textbf{数据生成机制}的假设。国债收益率是连续变化的慢变量,节假日缺口用时间线性插值无偏;而联邦基金利率是 FOMC 会议之间保持不变的政策阶梯,若同样插值,等于在两次会议之间虚构了从未发生的渐进加息——正确做法是前向填充。同理,月频与日频序列合并时,前向填充隐含``信息发布后保持有效直到更新''的假设,这与事件研究中的信息集设定是一致的。

\subsection{金融数据中``离群''不等于``错误''}
IQR 规则标记出的 21 个收益率极值日,逐一核对后全部是真实市场事件(2020-03 COVID、2022 加息周期等)。教科书式的``删除离群值''在这里恰恰是错的:BSM 失效分析(第 4 周)研究的正是这些日子。由此形成的操作规范:离群值检测的产出应当是\textbf{审计报告}(哪些点、为什么保留/修正)而非静默删除——决策要显式、可追溯。

\subsection{VIX 作为情绪代理的经济学依据}
经导师确认,情绪特征改由 VIX 派生(替代新闻数据)。这个替代在经济学上是有依据的:VIX 由标普 500 期权价格反解,包含风险中性方差预期与\textbf{方差风险溢价}两部分,后者正是市场为``对冲恐慌''支付的价格,所以 VIX 被称为恐惧指数并非修辞。构造上把 VIX 的原始水平(波动率角色)与标准化变换——滚动百分位(天然 0--1,即任务书的情绪分)、动量、状态标记(情绪角色)——分开使用,控制双重角色的信息重叠。这也呼应了第一部分 4.1 节的判断:期权业务的时序主线是波动率,``情绪''在可获得的数据里本质上仍是波动率信息的另一种切片;个股新闻能提供 VIX 之外的增量信息,但在免费数据约束下 VIX 是干净且理论自洽的选择。

\subsection{可复现管道的工程细节}
两个值得记录的设计:(1) 研究窗口与数据新鲜度解耦——建模固定使用 2018--2024,GitHub Actions 定时任务则通过环境变量把采集窗口延伸到当日,同一套代码同时服务回测与第 7 周的实时工具;(2) 模型可用集(1309 行零缺失)与全历史集(1760 行含预热期 NaN)分开输出,下游建模不必重复处理滚动窗口的边界问题。

\section{第 3 周:BSM Chooser 模型复现}

\subsection{``论文没给数字''时如何定义复现成功}
本周最先要解决的是方法论问题:参考论文没有报告一个明确的 chooser 价格,那么``与论文结果对照''对照什么?我采用的标准是\textbf{结构性复现}:(a) 逐行验证其 Table 3 路径表的计算逻辑;(b) 复现 Fig.3--6 四条敏感性曲线的方向与形状;(c) 在此之上用论文没有的\textbf{双重独立验证}(两种闭式解实现互查 + MC 收敛)确立自己实现的正确性。体会:复现的本质不是对齐数字,而是对齐\textbf{机制}——数字可以抄,机制对不对一验便知。

\subsection{正确性验证的``三角测量''}
本周实现的正确性由三条互相独立的证据链共同支撑:(1) Rubinstein 分解法与 Hull 直接公式两种解析实现差异 $7\times10^{-15}$(纯浮点误差);(2) 条件 MC 与终值 MC 两种模拟器都收敛到解析解 $\pm1$ 个标准误内;(3) 经济学排序 $\max(C,P)<\text{chooser}<\text{straddle}$ 及各敏感性方向全部符合无套利直觉。任何单一证据都可能因``实现与检验共享同一个错误''而失效,三角测量才可靠——这个习惯来自第一部分 6.1 节二叉树实验(收敛性 + 振荡结构双重检验)的延伸。

\subsection{条件蒙特卡洛:把解析知识注入模拟}
Chooser 的 MC 有个天然的降方差机会:$t_1$ 时被选中的那条腿的价值就是 BSM 公式,不必再模拟到 $T_2$。条件化后标准误降到终值模拟的 54\%(等效路径数 $\times3.4$),且完全无偏。这是初版 3.3 节``方差缩减是必需品''与 6.2 节控制变量实验的又一实例:\textbf{凡是问题中存在解析可解的子结构,就应把它从模拟中剥离}——模拟只留给真正解析不可解的部分。

\subsection{天真选择规则的经济学:一个状态依赖的结论}
论文在 $t_1$ 按 moneyness($S_{t_1}$ vs $K$)选 call/put,理性规则应比较两个期权的\textbf{价值},由 put--call parity 可得无差异点在 $S^*=Ke^{(q-r)\tau}$。用同种子配对(common random numbers)把 MC 噪声消到可分辨 0.02 量级后测得:论文参数($r=0.15\%$)下天真规则仅损失 6 bp——因为零利率时代 $S^*\approx K$;但 $r=10\%$ 时损失放大一个数量级到 60 bp。两点体会:(1) ``近似是否无害''是\textbf{状态依赖}的命题,对 2021 年数据无害的简化未必对 2022--2024 加息期无害;(2) 比较两个相近的量时,配对(CRN)比加大样本高效得多——这与统计里配对检验优于双样本检验是同一原理。

\subsection{对第一部分的呼应}
第一部分 3.2 节曾总结``有没有闭式解取决于收益依赖的路径统计量是否有可解分布''。chooser 是这个判断的一个干净正例:表面上它依赖``$t_1$ 的选择''这一路径信息,但 parity 把它压缩成 $S_{t_1}$ 的单时点函数,闭式解因此存在;而一旦允许在\textbf{任意}时刻选择(美式 chooser),该压缩失效,就回到初版疑问 1(Longstaff--Schwartz)的领域——enhancement 阶段的候选方向。

\section{第 4 周:基线模型评估}

\subsection{没有市场价格时,如何定义``误差''}
CME 成交价不可得后,本周要回答的方法论问题是:基线的 MAE/RMSE 到底衡量什么?最终采用``\textbf{事后真值(oracle)}''设计:同一份合约、同一个定价公式,唯一的区别是 $\sigma$ 用期权存续期内\textbf{实际实现}的波动率。这样构造的误差有干净的解释——它\textbf{只}度量常数波动率假设(用今天的回望波动率代替未来真实波动率)的代价,不混入公式本身或其他参数的误差。度量目标单一,后面 ML 的改进也就有明确的归因对象。

\subsection{一个被数据推翻的假设}
我预设``所有估计器在恐慌期误差更大'',并把它写成了审计断言——结果 252 日估计器把断言打断了:它在\textbf{平静期}误差更大(10.94 vs 8.55)。追查后确认不是 bug:252 日窗口在危机结束后一整年内仍拖着危机数据,于是在\textbf{危机后的平静期}持续高估;而快估计器(21日/EWMA)的失效集中在恐慌期(追高后市场均值回复)。两类失效方向相反,``没有单一窗口全状态占优''由此成为本周最重要的结论——它不是从文献里抄来的,是被自己的断言失败逼出来的。体会:\textbf{把预期写成会失败的断言,是比``看图确认''更诚实的自检};断言被推翻的时刻,要先怀疑代码,排除 bug 之后,剩下的就是发现。

\subsection{误差的可约与不可约部分}
2019 年尾段的巨大误差(MAE \$31.8)源于存续窗口内的 COVID 崩盘——任何回望方法都不可能预见。这提示基线误差应拆成两部分理解:\textbf{可预测的波动率动态}(聚集、均值回复、状态切换——ML 的合理战场)与\textbf{不可预见的跳变}(黑天鹅——任何时序模型的天花板)。这个拆分呼应了初版 4.3 节对 DL 的保留意见:对模型能力边界的预设,应该在评估设计里显式体现,而不是让一个混合的总指标掩盖``模型差''与``问题本身难''的区别。

\subsection{呼应与推进}
第 2 周埋下的两条线本周兑现:\texttt{fear\_regime} 特征成为失效分析的分组变量;``让模型学该听哪个窗口''的特征设计动机,现在有了定量依据(各窗口分状态的 MAE 表)。第 3 周的规则代价结论也在真实利率状态上复检成立(14.4bp $\to$ 32.2bp)。至此基线闭环完成,第 5--6 周 ML 的评估口径(同一 oracle、同一面板、同一指标)已经锁定,对比将是同台同尺的。

\section{第 5 周:机器学习增强定价(上)——目标参数化的教训}

\subsection{一次集体失败的正确打开方式}
初版七个模型(两条路线、含 LSTM)在验证集上\textbf{全部}输给 EWMA 基线,且失败模式一致(正偏差 $+14$ 至 $+26$ 美元)。处理顺序值得记录:第一步不是换模型、调参数,而是\textbf{跑穿六项审计排除实现错误}(数据泄漏、scaler 污染、切分越界、确定性等);全部通过后,失败才有资格被解读为``设计问题''。诊断结论:训练窗被 COVID 主导而验证窗是低波动状态,``直接预测波动率水平''的目标设计在状态漂移下注定失败——树模型结构上无法输出低于训练所见的值,这是模型类的固有性质,不是容量或调参问题。

\subsection{锚定:把已知结构从学习任务中剥离}
修正方案是把目标从水平改为对数修正量 $y=\ln(\sigma_{oracle}/\sigma_{EWMA})$。它的美妙之处在于\textbf{零点性质}:模型什么都不会时(输出 0)恰好还原 EWMA 基线,因此表现有了下界;模型学到的一切都是相对锚的\textbf{增量知识}。锚定 Ridge 由此成为首个击败基线的模型(综合分 0.567,定价 MAE $-39\%$,$R^2$ 从 $-1.87$ 到 $-0.07$)。这与第 3 周条件蒙特卡洛是同一个哲学的两次应用:\textbf{凡问题中存在已知/解析可解的子结构,就把它剥离,把模拟或学习只留给真正未知的残差}。事后看,残差学习(ResNet)、boosting、控制变量法都是这一原则的变体——它可能是本项目到目前为止最有迁移价值的一条经验。

\subsection{对初版疑问 4 的部分回应}
初版第 7 节疑问 4 问``Transformer/深度模型在低信噪比金融数据上是否优于简单模型''。本周给出了一个局部证据:在 664 个训练样本、强分布漂移的设定下,\textbf{被正确参数化的线性模型胜过所有非线性模型}(LSTM 2.45、树 3.7--4.2 vs 锚定 Ridge 0.567)——容量不是瓶颈,归纳偏置与目标设计才是。这不是终局(第 6 周滚动重训后树模型可能翻身),但它印证了初版 4.3 节的直觉:小样本、非平稳数据上,把结构性知识写进任务设计,比堆模型容量更有效。

\subsection{测试集纪律}
本周另一条自律:测试集(2024-03 至 2024-12)\textbf{封存不动},所有决策只用验证集;第 6 周调参定型后只打开一次。防的是``反复窥视测试集''造成的隐性过拟合——评估的可信度,取决于测试集被看过的次数。

\section{第 6 周:机器学习增强定价(下)——负结果的价值}

\subsection{验证集过拟合:亲历一次}
第 5 周锚定 Ridge 的验证集胜利(综合分 0.567)没有活过测试集:严格的 walk-forward 协议下,按验证表现选出的模型(LSTM,val MAE 2.82)在测试集上劣化到 17.07,且\textbf{没有任何配置}击败 EWMA 基线。事后诊断进一步显示,W5 的胜利部分是``训练分布恰好偏正''与验证期需求的巧合。这就是教科书里``对验证集过拟合''与``单一验证窗不构成证据''的活体样本——比在书上读到深刻得多。

\subsection{负结果为什么有价值}
三点:(1) 它经过了六项审计的资格审查,是``排除实现错误后的科学结论'',而非``没做出来'';(2) 它的归因是可操作的——锚定比率目标的低频状态自相关 + 只有一个半状态轮换的训练数据 + 标签窗口内的不可预见事件(2025-04 崩盘,与第 4 周 2019 尾段同构),三条各自指向明确的改进条件(更长历史、隐含波动率数据);(3) 它直接决定了工程决策:第 7--8 周工具以 EWMA 锚定 BSM 为主定价、ML 修正为带警示的参考——\textbf{知道模型在哪里不可信,和知道它在哪里可信同样是研究产出}。

\subsection{对初版疑问的回应(疑问 2/3/4)}
\textbf{疑问 4(简单 vs 深度模型)}:两周证据链完整了——W5 显示正确参数化的线性模型胜过一切非线性模型;W6 显示在严格样本外协议下\textbf{全体}被无参数的 EWMA 锚击败。在一个周期的低信噪比数据上,模型容量的边际价值为负,初版 4.3 节的怀疑得到强实证。\textbf{疑问 3(无套利约束)}:本项目的答案是结构性的——路线 A 把 ML 输出限制为波动率修正量,任何预测经 BSM 映射后自动满足无套利;这优于对端到端价格加软约束(路线 B 的锚定指数化只保证正价,不保证凸性)。\textbf{疑问 2(波动率微笑的一致性)}:受数据所限只有部分回应——VIX 是指数层面的隐含波动率代理,个股微笑数据不可得;这恰好是负结果归因中``改进条件''的一部分,留给 enhancement 阶段。

\subsection{SHAP 读出的规律与它的边界}
SHAP 显示模型学到的是干净的均值回复逻辑(当前波动率低于长期水平 $\to$ 预测正修正),与第 2 周互信息结论一致。规律方向正确、幅度不可标定——这提示失败不在特征或结构,而在\textbf{数据里可辨认的状态样本数}。样本外有效性的瓶颈可以不是算法,而是历史本身的多样性。

\section{第 7 周:敏感性分析与工具化}

\subsection{敏感性分析读出的合约“性格”}
极端情景把 chooser 的结构特征量化得很清楚:波动率 $+50\%$ 抬价近半,而 200bp 加息只动 1.5\%——call 与 put 两腿的 vega 同向叠加、rho 反向对消,所以 chooser 是一个近乎“纯波动率工具”;股价上下 15\% 都抬价(双向凸性),与第 3 周 $K$ 敏感性的 U 形互为印证。方法论上,把第 6 周 SHAP 对 $\ln\sigma$ 的贡献经定价函数映射成\textbf{美元}影响,比“重要性分数”更能与业务对话——锚定口径下长记忆波动率主导、VIX 让位,恰是锚定剥离已知结构的又一证据(与第 2 周水平口径的 VIX 主导不矛盾)。

\subsection{工具工程:UI 与逻辑分离的价值}
Streamlit 工具刻意做成“薄壳”:全部定价逻辑在 \texttt{src/tool/pricing\_service.py},UI 只调用服务函数。收益立现——六项审计(与闭式解逐位一致、无套利排序、情景重算、压力面单调、确定性、UI 零公式)全部可以脱离浏览器自动化运行;“UI 零公式”这条断言直接把“界面里偷偷写了第二套公式导致口径漂移”这类经典事故挡在门外。研究代码到产品代码的过渡,核心不是加界面,而是\textbf{让每一个数字仍然可审计}。

\section{第 8 周:交付与“1+1 融合”的回答}

\subsection{把八周收束成一个答案}
最终报告的写作让我把八周压缩成一条因果链:复现建立可信实现 $\to$ oracle 基线量化假设代价并揭示失效结构 $\to$ 封存测试集协议下 ML 未证优且归因清晰 $\to$ 工具据此选择 EWMA 锚定 BSM 为主定价。回看,每一步的\textbf{证据等级}都决定了下一步的设计——这大概是比任何单周结论都重要的收获:研究的推进单位不是“做了什么”,而是“确立了什么”。

\subsection{对导师“1+1”提示的回答}
导师的提示(“ML 与随机过程做 1+1 提效而非对立”)促使我把散落八周的线索整理成分工结构:\textbf{随机过程管框架}——无套利、可对冲、误差可归因;\textbf{ML 管残差}——框架覆盖不到的部分。五条落地路线按可行性排序:锚定残差学习(已实践,无套利自动满足、表现有下界)、收缩融合($\lambda$ 插值,最坏情形受控)、状态门控校准(ML 判状态、SP 换窗口——直击第 4 周“无单一窗口占优”)、隐含波动率曲面残差(信息与样本双优,需个股期权数据)、surrogate 加速(训练数据自生成,无信噪比问题)。值得注意的是,这五条与初版 4.3 节“DL 合理落点”的三个猜想(波动率曲面、定价加速、对冲学习)高度重合——八周的实证绕了一圈,验证并细化了项目开始前的直觉。

\subsection{关于负结果的最后一课}
导师对负结果的回应(“这是一定的,这代表你的探索结论和方向是对的”)本身就是一课:在量化研究的语境里,\textbf{被严格协议背书的负结果是方向性资产}——它告诉团队哪里不必再投入、以及投入什么(数据而非算法)才能改变结论。八周里最有价值的产出,可能不是那个 \$29.13,而是这套让正负结果同样可信的验证纪律。

\section{第 9 周:从研究原型到在线系统}

\subsection{一致性是迁移的唯一验收标准}
把数据从文件搬进 PostgreSQL、把报价放进 Redis,这件事本身没有任何研究含量,但它有一个极容易被忽略的失败方式:数字悄悄变了。本周的做法是把“迁移前后数字不变”写成可执行的断言——库中特征与文件特征在 2170 个交易日上最大差 $5.7\times10^{-14}$,库路径与文件路径的报价最大差 $2.3\times10^{-13}$,库中每条报价都能由同日特征重算复现。审计初版在末尾两个交易日失败,排查发现是 FRED 与 Yahoo 对近期数据的正常修订(流水线重拉了 10 天重叠区),于是“数据源有修订窗口”这个事实也进入了审计的定义。工程迁移的验收,和第 3 周复现闭式解、第 6 周封存测试集是同一种纪律:\textbf{每一个数字都要能被独立地重新算出来}。

\subsection{先测量,再优化}
导师对性能的要求是只看相对提升,这迫使我先留下一份“优化前”的基线。基线纠正了我的一个预判:我在代码注释里写过“每次报价重读模型文件是最大热点”,实测它每次只多 0.4 毫秒;真正相差三个数量级的是面板定价的标量循环与向量化闭式解(345 对 0.32 毫秒),以及每天全量重建特征表。如果没有测量就动手,我会花时间优化一个不存在的瓶颈。这和第 4 周“被数据推翻的假设”是同一课的工程版:直觉给出候选,测量给出顺序。

\subsection{理论匹配的锚 vs 经验锚:GARCH 期限结构}
初版报告 4.1 节记下过“GARCH 建模条件方差,波动率有均值回复结构”,本周第一次把它放到定价问题里检验。EWMA 是 GARCH(1,1) 的退化形式($\omega=0$,$\alpha+\beta=1$),没有均值回复;一年期 chooser 需要的是存续期内的平均方差,理论上匹配的锚是 GARCH 的期限结构预测——把当前方差沿 252 天向长期水平回复后取平均。在第 4 周的 oracle 面板上,这个锚的综合分 0.804,四种权重下均胜 EWMA。

但图形给出了更诚实的解释:持续性 0.95 在 252 天的平均下几乎被做满,1 年锚接近一条围绕长期水平的平线;它赢在\textbf{不追逐波动率尖峰}(P95AE 10.1 对 14.4),而不是赢在预测。代价是 $+3.43$ 美元的正偏差,2023 年这种持续平静期它反而比 EWMA 差。这与第 4 周“快估计器恐慌期追高”的发现是同一件事的两面,也再次提醒第 5 周的教训:单一标的的一次胜利不构成部署证据,第 10 周在 S\&P 500 面板复检之后,才决定是否换锚。

\section{第 10 周:从一只股票到五百只}

\subsection{单票看不到的东西}
把同一套代码从 JPM 扩到 590 只标的,代码几乎没有变,变的是数据。无分红公司的股息同比是 $0/0$,首次分红的公司是 $+\infty$,一条坏价格行让 Parkinson 波动率变成无穷,一家清算中的公司股息率 228\%。这些在 JPM 上永远不会出现,所以第 2 周的特征定义在 JPM 上是“对”的,却不是“完备”的。每一处修正都要过同一道关:JPM 的文件特征逐位不变(第 9 周审计 1)。这周我对“通用性”的理解具体了——不是模型能不能用在别的股票上,而是每一个定义在别的股票上是否仍然良定义。

\subsection{数据是有版本的}
第 9 周的两项审计本周失败,追查后发现 Yahoo 在每次分红后重算\textbf{整段}复权收盘价。这意味着本地文件快照和在线数据库是两个版本,跨版本比对特征根本不是在检验代码。修法是回到问题本身:审计要检验的是算术等价,那就在同一份价格上跑两套算术。顺带也解释了第 4 周存档审计为什么会失效,以及为什么历史报价行应当被当作“当日的审计留痕”而不是可以随时重算的事实。在此之前我把“确定性重跑”理解成“同一份代码跑两次一样”,现在加上一条:\textbf{同一份数据版本}。

\subsection{导师的“补充实验”与倍数的诚实算法}
五组优化前后对比都在同一进程里把“前”路径重测一遍,倍数与硬件无关,这是导师“只看相对提升”要求的技术含义。数字本身有起有落:情景表向量化 164 倍,全池报价却只有 5.4 倍——向量化定价不到 1 毫秒,412 毫秒几乎全在数据库扫 124 万行找每只标的的最新行。测量再一次纠正了直觉:瓶颈从计算换到了 I/O,下一步是索引而不是更多的 numpy。日更从外推的 44 分钟到实测 38 秒,兑现了对导师的承诺,但真正值得记的是这 38 秒里网络占 22 秒——余下的优化空间在并发而非算法。

\subsection{在 JPM 上赢的锚,在面板上没有赢}
第 9 周 GARCH 期限结构锚综合分 0.804,我在报告里写“单一标的一次胜利不构成部署证据”。本周 497 只标的给出了答案:胜率 54.5\%,中位 0.952,双峰分布,60 只长尾标的持续性接近 1 且全部正偏差。读初版报告 4.1 节时我记下“GARCH 有均值回复结构”,本周学到的是它的另一半:持续性接近 1 时长期方差水平估不准,而一年期限结构的平均会把这个错误的水平“做满”。反而不做期限结构平均的 1 日锚更稳。这个结果比“赢了”更有用——它给出了一个可检验的门控条件(持续性远离 1 才用期限结构锚),把一个“锚选哪个”的问题变成了一个“什么条件下选哪个”的问题。

\subsection{扩大截面不能替代更长的历史}
锚定 ML 第一轮在 494 只标的、9.5 万行验证集上全部劣于 EWMA,并且方向一致地高估。特征截尾只把 Ridge 从 7.16 改善到 6.12,说明主因不是横截面离群值,而是训练窗(2020--22)与验证窗(2023--24)之间的分布漂移——与第 5 周单票时是同一个机制。我原以为样本从 1 只到 494 只会解决第 6 周归因里的“状态样本数不足”,但 494 只股票在同一段时间里经历的是同一个宏观状态,横截面上的样本并不独立。这直接决定了第二轮的设计:学“相对同日其他标的”的修正而非绝对水平。

\section{第 11 周:测量决定优化,重述决定学习}

\subsection{被测量否定的优化也是结果}
本周四项优化里有一项失败了:把特征阶段的数据库往返从 1770 次减到 18 次,阶段时间纹丝不动。事后分解才看清,每票 24 毫秒里有 9 毫秒读、4 毫秒算、其余是数据库端执行 upsert 的成本——往返从来不是瓶颈。反倒是最“不像优化”的一项收益最大:写入时顺手维护一张“最新状态表”,让全池报价从扫 124 万行变成读 578 行,77 倍。这周我对“性能优化”的理解变成了两句话:先测,再改;改完必须再测,因为直觉对 I/O 的判断经常是错的。

\subsection{学什么,比用什么模型学更重要}
第 10 周第一轮锚定 ML 在 494 只标的上全军覆没,第二轮什么模型都没换,只改了学习目标:不学“该修多少”,只学“相对同一天其他标的该多修还是少修”。同一个 EWMA 锚下 Ridge 的综合分从 7.16 变成 0.859。这个改动的逻辑来自第 10 周的归因——失败源于时间上的分布漂移,而横截面去均值恰好把随时间漂移的那一部分从目标里减掉了。第 5 周“锚定 = 剥离已知结构”的哲学在这里又用了一次:锚剥离水平,去均值剥离市场共同因子,留给学习器的只剩真正的横截面差异。

\subsection{第一次赢了,以及为什么要把功劳分清楚}
测试集开封一次,GARCH 1 日锚加横截面 GBDT 的综合分 0.755,这是十一周里第一次有配置在封存测试集上击败 EWMA。但表格里紧挨着的一行是 GARCH 1 日锚单独 0.791:换锚拿走了大部分收益,ML 只贡献了 0.036。如果只报第一个数字,结论会是“机器学习赢了”;把两行放在一起,结论是“随机过程管水平,机器学习管横截面残差”——后者才是对的,也正是最终报告第 7 章写下的分工。测试期只有 7 个月,我把这条保留意见写在了结论旁边。复检时还发现一个更隐蔽的问题:测试期对 GARCH 锚并不是第一次见面——第 10 周提名它作候选时用的面板已经覆盖了这段时间,所以不能拿第 10 周的全样本结果当独立佐证;真正独立的证据是测试前的训练期(0.840)和验证期(0.777)。“测试集只开封一次”管住了模型选择,却没管住“候选是怎么来的”。

\subsection{最强的锚不能用}
用学校账号下的 OptionMetrics 数据检验了隐含波动率锚:JPM 上综合分 0.710,MAE 从 9.51 降到 5.90,按综合分是本项目检验过的锚里最强的(按 MAE 与同样平稳的 GARCH 1 年锚持平)。我起初把“它与 EWMA 该修多少的相关系数 0.78”当成它信息量的证据,复检时换了一个问题:如果隐含波动率毫无信息,这个相关会不会也这么大?答案是会——两边都含 EWMA 这个分母,把隐含波动率换成常数,相关系数反而是 0.83。真正的读数是:它赢在平稳而不在择时。三条序列的均值几乎相同,隐含波动率的标准差却只有 EWMA 的三分之一不到,与事后真值的时序相关只有 0.04。这是本周最有价值的一课:比值之间的相关,先查共享项;每写一句“解释了多少”,先问零假设下它会是多少。许可只允许本机研究,它不能进工具,只能作为方向性证据写进结项报告。

\section{第 12 周:结项——先登记,再检验,再相信}

\subsection{先写判据,再跑数}
导师对第 11 周提了两个问题:ML 增量的显著性边界,以及隐含波动率是否值得作为候选锚。这一周我做的第一件事不是跑数,而是把假设、统计量、推断方法、合格门槛和决策规则写进一个文件并记下它的哈希,之后才第一次运行分析。理由很朴素:如果先看到结果再决定“怎样算显著”,判据就会不自觉地向结果靠拢。事先写下来的规则后来真的约束了我:换锚和“隐含波动率对免费平稳锚”两项,估计值都很大、逐年为正,但它们的推断程序没有通过合格性检验,按规则只能作描述——如果没有事先的规则,我很可能会给它们贴上“显著”的标签。

\subsection{检验程序本身也要被检验}
记忆里有一条规则:自己写的统计程序先在已知答案的假数据上跑一遍。这一周它抓到了一个真实的错误:我用“行业 $\times$ 日期的均值”估计行业共同冲击,却忽略了均值里混着个股噪声,于是在纯噪声的数据上,检验的实际拒绝率明显高于名义的 5\%。修正之后再看真实数据,按股票重抽样的程序在四个假设上全部严重失准(拒绝率高达 \rnHTwoSizeTicker\% 一类的量级),只有“先抽行业再抽股票”的程序在两个假设上合格。如果直接用按股票的 bootstrap,我会报告一组漂亮但不可信的区间。

\subsection{“显著”的边界在哪里}
最终的回答是:2023--2024 年样本外,ML 相对 GARCH 锚每份期权少错 \rnHTwoDelta{} 美元,95\% 区间 \rnHTwoLo--\rnHTwoHi。但这是“以这段时期为条件”的结论——它说明在这两年里增量普遍存在于各只股票,却不能说明它会在别的年份重现:一年期标签下互不重叠的窗口只有两个,而且增量主要来自 2024 年。给出一个数,同时说清它回答了什么、没回答什么,我觉得这比一个没有边界的“显著”更接近导师要的东西。

\subsection{最强的锚为什么还不是候选}
隐含波动率在面板上是检验过的最强的锚,但分解显示,它的优势来自“每只股票的波动率水平定得准”,而不是“跟得上时间变化”——这也修正了我第 11 周在 JPM 一只股票上“赢在平稳”的读法。它最终被定位为研究基准:一是对免费平稳锚的优势没能通过合格的推断,二是许可与数据源使它无法部署。“知道最好的答案在哪里、同时清楚为什么现在不能用”,这件事在第 11 周还只是一句感想,这一周有了可以写进结论的证据。

\subsection{最难发现的错误都不在模型里}
这一周修掉的几个错误,没有一个是建模错误:调度器因为虚拟机随宿主睡眠而让单调时钟停走,一周没有运行也没有报警;写库时冲突更新会把没有提供的列清空;研究面板在数据库加列后出现列名冲突;我曾用“有没有行业字段”去证明“是不是当前成分股”,这是用被检验的东西证明它自己。它们的共同点是:单看每一步都“能跑”,只有换一个独立的角度去核对才会暴露。

\subsection{第一部分四个疑问的最终落点}
\begin{itemize}
  \item \textbf{疑问 1(美式期权的 MC 与 Longstaff--Schwartz)}:本项目没有做到,列为后续方向(结项报告第 13 节)。
  \item \textbf{疑问 2(波动率微笑与定价一致性)}:部分回应。第 11--12 周拿到个股期权的隐含波动率,发现它对“水平”的定位远好于任何基于历史的估计;但我只用了 365 天、50 Delta 的平值水平,没有研究微笑的形状,所以“给奇异期权定价时用哪个波动率”仍然只回答了一半。
  \item \textbf{疑问 3(DL 定价的无套利约束)}:第 5--6 周的结构性回答在第 12 周仍然成立——ML 只输出波动率修正量,任何预测经 BSM 映射后自动满足无套利;部署的横截面模型沿用这一结构。
  \item \textbf{疑问 4(简单模型与深度模型)}:十二周的证据一致:在低信噪比、状态样本少的数据上,模型容量不是瓶颈。真正带来改进的两步,一是换一个更好的随机过程锚,二是重新表述学习目标(只学横截面),都与模型复杂度无关。
\end{itemize}

\end{document}
